\documentclass[12pt]{article}

% Idioma e codificação
\usepackage[brazil]{babel}
\usepackage[utf8]{inputenc}
\usepackage[T1]{fontenc}

% Matemática e listas
\usepackage{amsmath, amssymb}
\usepackage{enumitem}

% Layout e cores
\usepackage{geometry}
\usepackage{graphicx}
\usepackage{xcolor}

% Diagramas e Portas Lógicas
\usepackage{tikz}
\usetikzlibrary{shapes.gates.logic.US, positioning, arrows.meta, calc, fit, shadows.blur}

% Cores institucionais
\definecolor{maincolor}{RGB}{0, 79, 151}
\definecolor{meuazul}{RGB}{0, 79, 151}

% Cabeçalho e rodapé
\usepackage{fancyhdr}

% Blocos
\usepackage[most]{tcolorbox}
\newtcolorbox{block}[1]{%
  colback=blue!5,
  colframe=maincolor,
  coltitle=white,
  title=#1,
  fonttitle=\bfseries,
  boxrule=1.5pt,
  arc=3mm,
  left=6pt,
  right=6pt,
  top=6pt,
  bottom=6pt
}

\newtcolorbox{resposta}[1]{%
  colback=white,
  colframe=meuazul!30,
  arc=1mm,
  height=#1,
  width=\textwidth,
  boxrule=0.8pt
}

% Margens
\geometry{margin=2.5cm}

% Fancyhdr
\pagestyle{fancy}
\fancyhf{}

\renewcommand{\headrulewidth}{2pt}
\renewcommand{\footrulewidth}{2pt}

\renewcommand{\headrule}{\hbox to\headwidth{%
  \color{meuazul}\leaders\hrule height \headrulewidth\hfill}}

\renewcommand{\footrule}{\hbox to\headwidth{%
  \color{meuazul}\leaders\hrule height \footrulewidth\hfill}}

\setlength{\headheight}{52pt}
\addtolength{\topmargin}{-20pt}

\fancyhead[L]{%
  \includegraphics[height=1.2cm]{./../logo.png}
}

\fancyhead[R]{%
  \raggedleft
  \textcolor{meuazul}{%
    \small
    \textbf{Lista de Exercícios — Desenvolvimento de Software para Dispositivos Móveis}\\
    Me.~Sergio Souza Novak\\
    \textit{Engenharia da Qualidade e Confiabilidade}
  }
}

\fancyfoot[C]{\textcolor{meuazul}{\thepage}}

% =========================
\begin{document}
% =========================

\vspace*{0.2cm}

\begin{center}
  {\Large \textbf{Noções de hardware de dispositivos móveis: processadores, memória, armazenamento, sensores e bateria. Sistemas operacionais móveis: Android e iOS, arquitetura e características}}\\[0.2cm]
  \textit{Instruções: Justifique suas respostas de forma clara e objetiva.
  Para questões discursivas, escreva na caixa de resposta abaixo de cada enunciado.}
\end{center}

\vspace{0.4cm}

\begin{center}
  \textit{``Boa atividade! Att. Me. Prof. Sergio Souza Novak.''}
\end{center}

% =========================

\textbf{Questão 1 — Sensores e Detecção de Metais (Múltipla Escolha)}
Uma equipe de desenvolvimento deseja criar um aplicativo para smartphones capaz de auxiliar
na detecção de metais utilizando os sensores já disponíveis no dispositivo móvel.
Durante o projeto, foram identificados os seguintes sensores:

\begin{itemize}[leftmargin=1.5em]
  \item Um sensor que mede a intensidade do campo magnético ao redor do celular.
  \item Um sensor que detecta movimentos e acelerações do aparelho.
  \item Um sensor que identifica a orientação e a rotação do dispositivo.
\end{itemize}

\vspace{0.4cm}

% ── Wireframes das telas do aplicativo ──────────────────────────────────────
\begin{center}
\begin{tikzpicture}[
  font=\sffamily\scriptsize,
  every node/.style={inner sep=0pt},
  phone/.style={draw=black!70, rounded corners=4mm, fill=white,
                line width=1.2pt, minimum width=3.8cm, minimum height=7cm},
  topbar/.style={draw=black!60, fill=black!10, line width=0.8pt,
                 minimum width=3.4cm, minimum height=0.55cm, rounded corners=1mm},
  widget/.style={draw=black!50, fill=white, line width=0.6pt,
                 minimum width=3.2cm, rounded corners=1mm},
  btn/.style={draw=black!70, fill=black!15, line width=0.8pt,
              minimum width=2.6cm, minimum height=0.55cm, rounded corners=2mm},
  sensorlbl/.style={draw=black!40, fill=black!5, line width=0.5pt,
                    minimum width=3.2cm, minimum height=0.45cm, rounded corners=1mm},
]

%% ─────── TELA 1: Escaneamento ───────
\node[phone] (phone1) at (0,0) {};

% câmera (notch)
\draw[fill=black!25, draw=black!50, line width=0.5pt]
  (0, 3.3) circle (0.12cm);

% barra de status
\node[topbar, anchor=north] (bar1) at (0, 3.1)
  {\textbf{MetalScan} \hfill 09:41};

% ícone de campo magnético (ondas)
\begin{scope}[shift={(0, 1.8)}]
  \draw[line width=1.2pt] (0,0) circle (0.38cm);
  \draw[line width=0.7pt] (0,0) circle (0.62cm);
  \draw[line width=0.4pt, dashed] (0,0) circle (0.86cm);
  \node[font=\sffamily\tiny\bfseries] at (0,-1.1) {Campo Magnético};
\end{scope}

% barra de intensidade
\node[widget, minimum height=0.38cm, anchor=north] (bar_bg) at (0, 0.55)
  {};
\node[draw=none, fill=black!30, minimum width=1.1cm, minimum height=0.28cm,
      anchor=west, rounded corners=0.5mm] at ([xshift=0.05cm]bar_bg.west) {};
\node[font=\sffamily\tiny, anchor=east] at ([xshift=-0.08cm]bar_bg.east)
  {intensidade};

% leituras dos sensores
\node[sensorlbl, anchor=north] (s1) at (0, 0.1)
  {\textbf{Magnetômetro:} \ 48 µT};
\node[sensorlbl, anchor=north] (s2b) at (0, -0.45)
  {\textbf{Acelerômetro:} \ 0.3 m/s²};
\node[sensorlbl, anchor=north] (s3) at (0, -0.98)
  {\textbf{Giroscópio:} \ 2.1 °/s};

% botão iniciar
\node[btn, anchor=north] at (0, -1.6)
  {\textbf{\ Iniciar Varredura}};

% legenda
\node[font=\sffamily\tiny\itshape, text=black!60, anchor=north] at (0,-3.0)
  {Tela 1 — Escaneamento};

%% ─────── TELA 2: Metal Detectado ───────
\node[phone] (phone2) at (5.5,0) {};

% câmera
\draw[fill=black!25, draw=black!50, line width=0.5pt]
  (5.5, 3.3) circle (0.12cm);

% barra de status
\node[topbar, anchor=north] (bar2) at (5.5, 3.1)
  {\textbf{MetalScan} \hfill 09:42};

% ícone de alerta
\begin{scope}[shift={(5.5, 1.75)}]
  \draw[line width=1.5pt, fill=black!8]
    (0, 0.55) -- (-0.55,-0.45) -- (0.55,-0.45) -- cycle;
  \node[font=\bfseries\small] at (0, 0.04) {!};
  \node[font=\sffamily\tiny\bfseries] at (0,-0.85)
    {METAL DETECTADO};
\end{scope}

% barra de intensidade alta
\node[widget, minimum height=0.38cm, anchor=north] (bar_bg2) at (5.5, 0.55) {};
\node[draw=none, fill=black!60, minimum width=2.8cm, minimum height=0.28cm,
      anchor=west, rounded corners=0.5mm]
  at ([xshift=0.05cm]bar_bg2.west) {};
\node[font=\sffamily\tiny, anchor=east, text=white]
  at ([xshift=-0.12cm]bar_bg2.east) {alta};

% leituras — estado de alerta
\node[sensorlbl, anchor=north] at (5.5, 0.1)
  {\textbf{Magnetômetro:} \ 312 µT};
\node[sensorlbl, anchor=north] at (5.5, -0.45)
  {\textbf{Acelerômetro:} \ 1.9 m/s²};
\node[sensorlbl, anchor=north] at (5.5, -0.98)
  {\textbf{Giroscópio:} \ 8.4 °/s};

% botão nova varredura
\node[btn, anchor=north] at (5.5, -1.6)
  {\textbf{\ Nova Varredura}};

% legenda
\node[font=\sffamily\tiny\itshape, text=black!60, anchor=north] at (5.5,-3.0)
  {Tela 2 — Metal Detectado};

%% seta de transição
\draw[-{Stealth[length=6pt]}, line width=1pt, black!50]
  ([xshift=0.25cm]phone1.east) -- ([xshift=-0.25cm]phone2.west)
  node[midway, above, font=\sffamily\tiny, text=black!50] {detecção};

\end{tikzpicture}
\end{center}
% ────────────────────────────────────────────────────────────────────────────

\vspace{0.3cm}
Considerando o funcionamento desse aplicativo para detecção de metais, assinale a alternativa que indica corretamente qual é o principal sensor responsável por medir as variações e identificar a presença do metal ao redor do aparelho:

\vspace{0.4cm}

\begin{enumerate}[label=\alph*)]
    \item Acelerômetro
    \item Magnetômetro
    \item Giroscópio
    \item Sensor de Proximidade
    \item Barômetro
\end{enumerate}

% =========================

\vspace{0.8cm}
\noindent\rule{\textwidth}{0.8pt}
\vspace{0.3cm}

% ── Questão 2 ───────────────────────────────────────────────────────────────
\noindent\textbf{Questão 2 — Nível de Bolha Digital (Discursiva)}

\vspace{0.3cm}
\noindent
Uma empresa de marcenaria deseja desenvolver um aplicativo para smartphones que auxilie os
profissionais na instalação de prateleiras, armários e bancadas. O aplicativo deverá informar
se uma superfície está perfeitamente nivelada, funcionando de maneira semelhante a um
\textit{nível de bolha} utilizado na construção civil.

\vspace{0.3cm}
\noindent
Considerando os sensores disponíveis em um smartphone, \textbf{indique qual é o principal
sensor que deve ser utilizado} para essa funcionalidade e explique por que ele é o mais
adequado para a aplicação proposta. Descreva também como os dados fornecidos por esse sensor
podem ajudar o marceneiro durante o trabalho.

\vspace{0.5cm}

% ── Wireframes: Nível de Bolha ───────────────────────────────────────────────
\begin{center}
\begin{tikzpicture}[
  font=\sffamily\scriptsize,
  every node/.style={inner sep=0pt},
  phone/.style={draw=black!70, rounded corners=4mm, fill=white,
                line width=1.2pt, minimum width=3.8cm, minimum height=7cm},
  topbar/.style={draw=black!60, fill=black!10, line width=0.8pt,
                 minimum width=3.4cm, minimum height=0.55cm, rounded corners=1mm},
  btn/.style={draw=black!70, fill=black!15, line width=0.8pt,
              minimum width=2.8cm, minimum height=0.55cm, rounded corners=2mm},
  infolbl/.style={draw=black!40, fill=black!5, line width=0.5pt,
                  minimum width=3.2cm, minimum height=0.45cm, rounded corners=1mm},
]

%% ─────── TELA A: Superfície Nivelada ───────
\node[phone] (phoneA) at (0,0) {};

\draw[fill=black!25, draw=black!50, line width=0.5pt] (0,3.3) circle (0.12cm);

\node[topbar, anchor=north] at (0, 3.1) {\textbf{NívelPro} \hfill 14:22};

% Círculo do nível de bolha — nivelado
\begin{scope}[shift={(0, 1.5)}]
  % área do visor
  \draw[draw=black!60, fill=black!5, line width=1pt] (0,0) circle (1.0cm);
  % linhas de referência
  \draw[black!30, line width=0.5pt] (-1.0,0) -- (1.0,0);
  \draw[black!30, line width=0.5pt] (0,-1.0) -- (0,1.0);
  \draw[black!20, dashed, line width=0.4pt] (0,0) circle (0.25cm);
  % bolha no centro (nivelado)
  \draw[draw=black!70, fill=black!20, line width=0.8pt] (0,0) circle (0.22cm);
  % rótulo
  \node[font=\sffamily\tiny\bfseries] at (0,-1.3) {NIVELADO};
\end{scope}

% Leitura do acelerômetro
\node[infolbl, anchor=north] at (0, 0.1)
  {\textbf{Acel. X:} \ 0.01 m/s²};
\node[infolbl, anchor=north] at (0, -0.43)
  {\textbf{Acel. Y:} \ 0.00 m/s²};
\node[infolbl, anchor=north] at (0, -0.95)
  {\textbf{Acel. Z:} \ 9.81 m/s²};

\node[btn, anchor=north] at (0, -1.6) {\textbf{Calibrar}};

\node[font=\sffamily\tiny\itshape, text=black!60, anchor=north] at (0,-3.0)
  {Tela A — Superfície Nivelada};

%% ─────── TELA B: Superfície Inclinada ───────
\node[phone] (phoneB) at (5.5,0) {};

\draw[fill=black!25, draw=black!50, line width=0.5pt] (5.5,3.3) circle (0.12cm);

\node[topbar, anchor=north] at (5.5, 3.1) {\textbf{NívelPro} \hfill 14:23};

% Círculo do nível de bolha — inclinado
\begin{scope}[shift={(5.5, 1.5)}]
  \draw[draw=black!60, fill=black!5, line width=1pt] (0,0) circle (1.0cm);
  \draw[black!30, line width=0.5pt] (-1.0,0) -- (1.0,0);
  \draw[black!30, line width=0.5pt] (0,-1.0) -- (0,1.0);
  \draw[black!20, dashed, line width=0.4pt] (0,0) circle (0.25cm);
  % bolha deslocada (inclinado)
  \draw[draw=black!70, fill=black!40, line width=0.8pt] (0.52,0.38) circle (0.22cm);
  % seta indicando correção
  \draw[-{Stealth[length=4pt]}, black!60, line width=0.6pt]
    (0.52,0.38) -- (0.08,0.06);
  \node[font=\sffamily\tiny\bfseries] at (0,-1.3) {INCLINADO};
\end{scope}

% Leitura do acelerômetro — inclinado
\node[infolbl, anchor=north] at (5.5, 0.1)
  {\textbf{Acel. X:} \ 2.43 m/s²};
\node[infolbl, anchor=north] at (5.5, -0.43)
  {\textbf{Acel. Y:} \ 1.87 m/s²};
\node[infolbl, anchor=north] at (5.5, -0.95)
  {\textbf{Acel. Z:} \ 9.45 m/s²};

\node[btn, anchor=north] at (5.5, -1.6) {\textbf{Ajustar Inclinação}};

\node[font=\sffamily\tiny\itshape, text=black!60, anchor=north] at (5.5,-3.0)
  {Tela B — Superfície Inclinada};

%% seta de transição
\draw[-{Stealth[length=6pt]}, line width=1pt, black!50]
  ([xshift=0.25cm]phoneA.east) -- ([xshift=-0.25cm]phoneB.west)
  node[midway, above, font=\sffamily\tiny, text=black!50] {inclinar};

\end{tikzpicture}
\end{center}
% ────────────────────────────────────────────────────────────────────────────

\vspace{0.4cm}

\noindent\textit{Resposta:}

\begin{resposta}{12cm}
\end{resposta}

% =========================

\vspace{0.8cm}
\noindent\rule{\textwidth}{0.8pt}
\vspace{0.3cm}

% ── Questão 3 ───────────────────────────────────────────────────────────────
\noindent\textbf{Questão 3 — Identificação de Pragas com IA e Geolocalização (Discursiva)}

\vspace{0.3cm}
\noindent
Uma empresa do setor agrícola pretende desenvolver um aplicativo para smartphones capaz de
identificar pragas em plantações por meio de Inteligência Artificial e Visão Computacional.
O usuário deverá apontar a câmera do celular para a planta, e o aplicativo analisará a
imagem para detectar possíveis pragas e apresentar recomendações de manejo. Além disso, o
sistema deverá registrar a localização geográfica de cada detecção para construir um
\textbf{mapa de calor} da infestação de pragas em toda a fazenda, auxiliando na tomada de
decisões sobre o controle da lavoura.

\vspace{0.3cm}
\noindent
Considerando os recursos de hardware disponíveis em um smartphone, \textbf{indique quais são
os principais componentes necessários} para a captura dos dados utilizados pelo sistema.
Explique a função de cada componente e descreva como eles contribuem tanto para a
\textbf{identificação das pragas} quanto para a \textbf{geração do mapa de calor} da
infestação na propriedade rural.

\vspace{0.5cm}

% ── Wireframes: App de Pragas ─────────────────────────────────────────────────
\begin{center}
\begin{tikzpicture}[
  font=\sffamily\scriptsize,
  every node/.style={inner sep=0pt},
  phone/.style={draw=black!70, rounded corners=4mm, fill=white,
                line width=1.2pt, minimum width=3.6cm, minimum height=6.8cm},
  topbar/.style={draw=black!60, fill=black!10, line width=0.8pt,
                 minimum width=3.2cm, minimum height=0.52cm, rounded corners=1mm},
  btn/.style={draw=black!70, fill=black!15, line width=0.8pt,
              minimum width=2.8cm, minimum height=0.52cm, rounded corners=2mm},
  lbl/.style={draw=black!40, fill=black!5, line width=0.5pt,
              minimum width=3.2cm, minimum height=0.42cm, rounded corners=1mm},
  chip/.style={draw=black!50, fill=black!8, line width=0.6pt,
               minimum width=1.4cm, minimum height=0.42cm, rounded corners=1mm},
]

%% ─── TELA 1: Captura ───
\node[phone] (p1) at (0,0) {};
\draw[fill=black!25, draw=black!50, line width=0.4pt] (0,3.2) circle (0.11cm);
\node[topbar, anchor=north] at (0,3.05) {\textbf{AgroScan} \hfill};

% visor da câmera
\draw[draw=black!50, fill=black!8, line width=0.8pt, rounded corners=1mm]
  (-1.5, 1.9) rectangle (1.5, 0.3);
% mira central
\draw[black!40, line width=0.5pt] (-0.25,1.1) -- (0.25,1.1);
\draw[black!40, line width=0.5pt] (0, 0.85) -- (0, 1.35);
\draw[black!30, dashed, line width=0.4pt]
  (-0.6,1.4) -- (-0.6,1.7) -- (-0.3,1.7);
\draw[black!30, dashed, line width=0.4pt]
  (0.6,1.4) -- (0.6,1.7) -- (0.3,1.7);
\draw[black!30, dashed, line width=0.4pt]
  (-0.6,0.8) -- (-0.6,0.5) -- (-0.3,0.5);
\draw[black!30, dashed, line width=0.4pt]
  (0.6,0.8) -- (0.6,0.5) -- (0.3,0.5);
\node[font=\sffamily\tiny, text=black!50] at (0,1.1) {aponte para a planta};

% GPS chip
\node[chip] (gps1) at (-0.82, -0.05) {\tiny GPS};
\node[chip] at (0.82, -0.05) {\tiny Câmera};

\node[lbl, anchor=north] at (0,-0.35)
  {\textbf{Localização:} \ -18.92°, -49.71°};

\node[btn, anchor=north] at (0,-0.88) {\textbf{\ Capturar e Analisar}};

\node[font=\sffamily\tiny, text=black!60, anchor=north] at (0,-2.15)
  {Tela 1 — Captura};

%% ─── TELA 2: Resultado da IA ───
\node[phone] (p2) at (5.3,0) {};
\draw[fill=black!25, draw=black!50, line width=0.4pt] (5.3,3.2) circle (0.11cm);
\node[topbar, anchor=north] at (5.3,3.05) {\textbf{AgroScan} \hfill};

% miniatura da imagem capturada
\draw[draw=black!50, fill=black!12, line width=0.7pt, rounded corners=1mm]
  (3.85,2.55) rectangle (6.75,1.35);
\node[font=\sffamily\tiny, text=black!50] at (5.3,1.95) {imagem capturada};
% bounding box da praga
\draw[draw=black!70, line width=0.8pt, dashed]
  (4.5,2.38) rectangle (5.45,1.52);
\node[font=\sffamily\tiny\bfseries] at (4.97,2.46) {praga};

% resultado
\node[lbl, anchor=north] at (5.3,1.45)
  {\textbf{Praga:} \ Lagarta-do-cartucho};
\node[lbl, anchor=north] at (5.3,0.95)
  {\textbf{Confiança:} \ 94\%};
\node[lbl, anchor=north] at (5.3,0.45)
  {\textbf{GPS:} \ -18.92°, -49.71°};

% recomendação
\draw[draw=black!40, fill=black!3, line width=0.5pt, rounded corners=1mm]
  (3.8,-0.08) rectangle (6.8,-0.75);
\node[font=\sffamily\tiny, text width=2.7cm, align=center] at (5.3,-0.42)
  {Recomendação: aplicar inseticida Bt em 48h};

\node[btn, anchor=north] at (5.3,-0.9) {\textbf{Ver no Mapa}};

\node[font=\sffamily\tiny, text=black!60, anchor=north] at (5.3,-2.15)
  {Tela 2 — Resultado IA};

%% ─── TELA 3: Mapa de Calor ───
\node[phone] (p3) at (10.6,0) {};
\draw[fill=black!25, draw=black!50, line width=0.4pt] (10.6,3.2) circle (0.11cm);
\node[topbar, anchor=north] at (10.6,3.05) {\textbf{AgroScan — Mapa} \hfill };

% grade do mapa (talhão)
\draw[draw=black!40, fill=black!5, line width=0.6pt, rounded corners=1mm]
  (9.2,2.55) rectangle (12.0,0.2);
% células de calor (intensidade simulada com fill)
\fill[black!50, rounded corners=0.5mm] (10.3,2.1) rectangle (11.1,1.5);
\fill[black!30, rounded corners=0.5mm] (9.5,1.7) rectangle (10.1,1.2);
\fill[black!20, rounded corners=0.5mm] (11.1,1.3) rectangle (11.7,0.9);
\fill[black!10, rounded corners=0.5mm] (9.6,0.9) rectangle (10.2,0.55);
\fill[black!10, rounded corners=0.5mm] (11.2,2.1) rectangle (11.7,1.7);
% pontos de detecção
\fill[black!80] (10.7,1.8) circle (0.09cm);
\fill[black!80] (9.8,1.45) circle (0.07cm);
\fill[black!80] (11.4,1.1) circle (0.07cm);
\fill[black!60] (9.9,0.72) circle (0.06cm);
% legenda do mapa
\node[font=\sffamily\tiny, text=black!50] at (10.6,0.1)
  {$\blacksquare$ alta \quad $\blacksquare$ média \quad $\square$ baixa};

% legenda inferior
\node[lbl, anchor=north] at (10.6,-0.3)
  {\textbf{Detecções:} \ 4 \ pontos registrados};
\node[btn, anchor=north] at (10.6,-0.82) {\textbf{Exportar Relatório}};

\node[font=\sffamily\tiny, text=black!60, anchor=north] at (10.6,-2.15)
  {Tela 3 — Mapa de Calor};

%% setas de transição
\draw[-{Stealth[length=5pt]}, line width=0.9pt, black!45]
  ([xshift=0.2cm]p1.east) -- ([xshift=-0.2cm]p2.west)
  node[midway, above, font=\sffamily\tiny, text=black!45] {analisar};
\draw[-{Stealth[length=5pt]}, line width=0.9pt, black!45]
  ([xshift=0.2cm]p2.east) -- ([xshift=-0.2cm]p3.west)
  node[midway, above, font=\sffamily\tiny, text=black!45] {mapear};

\end{tikzpicture}
\end{center}
% ────────────────────────────────────────────────────────────────────────────

\vspace{0.4cm}

\noindent\textit{Resposta:}

\begin{resposta}{12cm}
\end{resposta}

% =========================

\vspace{0.8cm}
\noindent\rule{\textwidth}{0.8pt}
\vspace{0.3cm}

% ── Questão 4 ───────────────────────────────────────────────────────────────
\noindent\textbf{Questão 4 — Aplicativo de Emergência: Palavra-Chave, Agitação e Geolocalização (Discursiva)}

\vspace{0.3cm}
\noindent
Uma organização de proteção às mulheres deseja desenvolver um aplicativo para smartphones
que permita solicitar ajuda em situações de emergência de forma \textbf{rápida e discreta}.
Para evitar acionamentos indevidos, o aplicativo só deverá enviar um alerta à polícia quando
a usuária pronunciar uma \textbf{palavra-chave} previamente cadastrada (por exemplo,
\textit{``estetoscópio''}) e, em seguida, \textbf{chacoalhar o celular}. Após a confirmação,
o sistema deverá enviar a \textbf{localização da vítima} para as autoridades e para contatos
de confiança.

\vspace{0.3cm}
\noindent
Considerando os recursos de hardware disponíveis em um smartphone,
\textbf{identifique os principais componentes e sensores} necessários para implementar essa
solução. Explique a função de cada elemento de hardware e descreva como eles atuam no
\textbf{reconhecimento da palavra-chave}, na \textbf{detecção do movimento de chacoalhar} o
aparelho e no \textbf{envio da localização} da usuária. Além disso, justifique por que cada
componente é importante para garantir o funcionamento adequado do aplicativo em uma situação
de emergência.

\vspace{0.5cm}

% ── Wireframes: App de Emergência ────────────────────────────────────────────
\begin{center}
\begin{tikzpicture}[
  font=\sffamily\scriptsize,
  every node/.style={inner sep=0pt},
  phone/.style={draw=black!70, rounded corners=4mm, fill=white,
                line width=1.2pt, minimum width=3.6cm, minimum height=6.8cm},
  topbar/.style={draw=black!60, fill=black!10, line width=0.8pt,
                 minimum width=3.2cm, minimum height=0.52cm, rounded corners=1mm},
  btn/.style={draw=black!70, fill=black!15, line width=0.8pt,
              minimum width=2.8cm, minimum height=0.52cm, rounded corners=2mm},
  btnalert/.style={draw=black!80, fill=black!35, line width=1pt,
                   minimum width=2.8cm, minimum height=0.52cm, rounded corners=2mm},
  lbl/.style={draw=black!40, fill=black!5, line width=0.5pt,
              minimum width=3.2cm, minimum height=0.42cm, rounded corners=1mm},
  chip/.style={draw=black!50, fill=black!8, line width=0.5pt,
               minimum width=1.4cm, minimum height=0.40cm, rounded corners=1mm},
]

%% ─── TELA 1: Modo Espera (Discreta) ───
\node[phone] (p1) at (0,0) {};
\draw[fill=black!25, draw=black!50, line width=0.4pt] (0,3.2) circle (0.11cm);
\node[topbar, anchor=north] at (0,3.05) {\textbf{SafeWord} \hfill $\blacksquare$};

% ícone de microfone (escuta ativa)
\begin{scope}[shift={(0, 1.7)}]
  \draw[draw=black!60, fill=black!8, line width=0.8pt, rounded corners=1mm]
    (-0.22,-0.42) rectangle (0.22,0.42);
  \draw[draw=black!60, line width=0.8pt]
    (-0.38,0.1) arc[start angle=180, end angle=0, radius=0.38cm];
  \draw[draw=black!60, line width=0.8pt] (0,0.48) -- (0,0.72);
  \draw[draw=black!60, line width=0.8pt] (-0.22,0.72) -- (0.22,0.72);
  % ondas de áudio
  \draw[black!30, line width=0.4pt]
    (-0.65, 0.0) arc[start angle=150, end angle=210, radius=0.55cm];
  \draw[black!30, line width=0.4pt]
    (0.65, 0.0) arc[start angle=30, end angle=-30, radius=0.55cm];
  \node[font=\sffamily\tiny, text=black!50] at (0,-0.7) {aguardando palavra-chave\ldots};
\end{scope}

% chips de hardware ativos
\node[chip] at (-0.82,0.38) {\tiny Microfone};
\node[chip] at (0.82,0.38)  {\tiny CPU};

\node[lbl, anchor=north] at (0,0.08)
  {\textbf{Status:} \ escutando\ldots};
\node[lbl, anchor=north] at (0,-0.42)
  {\textbf{Palavra-chave:} \ cadastrada};
\node[lbl, anchor=north] at (0,-0.94)
  {\textbf{GPS:} \ ativo em 2.º plano};

\node[btn, anchor=north] at (0,-1.55) {\textbf{Configurações}};

\node[font=\sffamily\tiny\itshape, text=black!60, anchor=north] at (0,-2.15)
  {Tela 1 — Modo Espera Discreta};

%% ─── TELA 2: Palavra-Chave Detectada → Agitar ───
\node[phone] (p2) at (5.3,0) {};
\draw[fill=black!25, draw=black!50, line width=0.4pt] (5.3,3.2) circle (0.11cm);
\node[topbar, anchor=north] at (5.3,3.05) {\textbf{SafeWord} \hfill $\blacksquare$};

% ícone de agitação (setas curvas)
\begin{scope}[shift={(5.3,1.65)}]
  % celular inclinado (retângulo rotacionado)
  \begin{scope}[rotate=18]
    \draw[draw=black!70, fill=black!10, line width=0.9pt, rounded corners=1.5mm]
      (-0.18,-0.45) rectangle (0.18,0.45);
  \end{scope}
  % setas de movimento
  \draw[-{Stealth[length=4pt]}, black!55, line width=0.7pt]
    (-0.55, 0.3) arc[start angle=140, end angle=40, radius=0.55cm];
  \draw[-{Stealth[length=4pt]}, black!55, line width=0.7pt]
    (0.55,-0.3) arc[start angle=-40, end angle=-140, radius=0.55cm];
  \node[font=\sffamily\tiny\bfseries] at (0,-0.8) {AGITE O CELULAR};
\end{scope}

% barra de progresso do shake
\draw[draw=black!40, fill=black!8, line width=0.5pt, rounded corners=0.5mm]
  (3.9, 0.42) rectangle (6.7, 0.16);
\draw[draw=none, fill=black!45, rounded corners=0.5mm]
  (3.9, 0.42) rectangle (5.8, 0.16);
\node[font=\sffamily\tiny, text=black!50] at (5.3,0.55) {intensidade do movimento};

\node[chip] at (4.48,0.0) {\tiny Acelerôm.};
\node[chip] at (6.12,0.0) {\tiny Giroscópio};

\node[lbl, anchor=north] at (5.3,-0.35)
  {\textbf{Palavra detectada!}  Confirme};

% contador regressivo
\draw[draw=black!50, fill=black!5, line width=0.7pt] (5.3,-0.92) circle (0.32cm);
\node[font=\sffamily\small\bfseries] at (5.3,-0.92) {5};
\node[font=\sffamily\tiny, text=black!50] at (5.3,-1.35) {segundos para confirmar};

\node[btnalert, anchor=north] at (5.3,-1.65)
  {\textcolor{white}{\textbf{CANCELAR}}};

\node[font=\sffamily\tiny\itshape, text=black!60, anchor=north] at (5.3,-2.15)
  {Tela 2 — Confirmação por Agitação};

%% ─── TELA 3: Alerta Enviado ───
\node[phone] (p3) at (10.6,0) {};
\draw[fill=black!25, draw=black!50, line width=0.4pt] (10.6,3.2) circle (0.11cm);
\node[topbar, anchor=north] at (10.6,3.05) {\textbf{SafeWord} \hfill $\blacksquare$};

% ícone de alerta enviado (escudo com check)
\begin{scope}[shift={(10.6,1.7)}]
  \draw[draw=black!65, fill=black!10, line width=1.1pt]
    (0,0.55) -- (-0.5,0.25) -- (-0.5,-0.3) -- (0,-0.55) -- (0.5,-0.3) -- (0.5,0.25) -- cycle;
  \draw[black!65, line width=1.2pt] (-0.18,0.0) -- (-0.02,-0.22) -- (0.22,0.28);
  \node[font=\sffamily\tiny\bfseries] at (0,-0.82) {ALERTA ENVIADO};
\end{scope}

\node[lbl, anchor=north] at (10.6,0.08)
  {\textbf{Polícia:} \ notificada};
\node[lbl, anchor=north] at (10.6,-0.38)
  {\textbf{Contatos:} \ 3 avisados};
\node[lbl, anchor=north] at (10.6,-0.88)
  {\textbf{Localização:} \ -18.91°, -49.72°};
\node[lbl, anchor=north] at (10.6,-1.36)
  {\textbf{GPS:} \ enviado às autoridades};

\node[btn, anchor=north] at (10.6,-1.86) {\textbf{Cancelar Alerta}};

\node[font=\sffamily\tiny\itshape, text=black!60, anchor=north] at (10.6,-2.5)
  {Tela 3 — Alerta Enviado};

%% setas de transição
\draw[-{Stealth[length=5pt]}, line width=0.9pt, black!45]
  ([xshift=0.2cm]p1.east) -- ([xshift=-0.2cm]p2.west)
  node[midway, above, font=\sffamily\tiny, text=black!45] {palavra-chave};
\draw[-{Stealth[length=5pt]}, line width=0.9pt, black!45]
  ([xshift=0.2cm]p2.east) -- ([xshift=-0.2cm]p3.west)
  node[midway, above, font=\sffamily\tiny, text=black!45] {agitar};

\end{tikzpicture}
\end{center}
% ────────────────────────────────────────────────────────────────────────────

\vspace{0.4cm}

\noindent\textit{Resposta:}

\begin{resposta}{12cm}
\end{resposta}

% =========================
\end{document}
% =========================
